% !TeX root = ../../Book2.tex
% 纲要标题：### 18.2 循环代数
% 纲要位置：计划纲要.md，第 1637 行。

\section{循环代数}
\label{b2:sec:1802}

一个循环 Galois 扩张不仅能提供标量，还能提供乘法的扭转规则。把扩张的生成自同构写成一个新元素的共轭作用，再规定该元素的某次幂，就得到循环代数。它是否分裂，最后化为原扩张中的一个范数方程。

% 纲要要点（供目录核对）：
%
% * 循环 Galois 扩张及扭曲乘法；先由生成关系产生进位公式，再以两次进位恒等式证明结合律，回扣代数展示的泛性质
% * 循环代数的中心单性；分别核对中心与最少非零系数证明，分裂映射使用原代数的右子域空间
% * 范数与分裂条件；一维子域空间的半线性算子产生范数方程，换生成元直接计算参数及恢复原生成元
%

\subsection{循环 Galois 扩张及扭曲乘法}
\label{b2:subsec:1802:01}

设 \(K/k\) 为次数 \(n\ge2\) 的循环 Galois 扩张，\(\operatorname{Gal}(K/k)=\langle\sigma\rangle\)，并取 \(a\in k^\times\)。希望在包含 \(K\) 的代数中加入一个元素 \(u\)，要求 \(ub=\sigma(b)u\)（\(b\in K\)）及 \(u^n=a\)。前一关系反复使用便给出 \(u^ic=\sigma^i(c)u^i\)；若 \(i+j=qn+r\)、\(0\le r<n\)，后一关系将 \(u^{i+j}\) 化成 \(a^qu^r\)。这两个规则决定了乘积的系数和指数。据此，在左 \(K\) 向量空间
\[
C=K\oplus Ku\oplus\cdots\oplus Ku^{n-1}
\]
上，用下式定义乘法，并按 \(k\) 双线性扩张：
\begin{equation}\label{b2:eq:cyclic-algebra-product}
(bu^i)(cu^j)=b\sigma^i(c)a^{\left\lfloor(i+j)/n\right\rfloor}u^{\,i+j-n\left\lfloor(i+j)/n\right\rfloor},
\qquad 0\le i,j<n.
\end{equation}
这里 \(q=\lfloor(i+j)/n\rfloor\) 记录进位次数，每次进位产生一个 \(a\)。乘法一般不是 \(K\) 双线性的，但 \(\sigma\) 固定 \(k\)，所以仍为 \(k\) 双线性。

\begin{proposition}\label{b2:prop:cyclic-algebra-construction}
设 \(K/k\) 为次数 \(n\geq2\) 的循环 Galois 扩张，\(\operatorname{Gal}(K/k)=\langle\sigma\rangle\)，\(a\in k^\times\)。在左 \(K\) 空间 \(C=\bigoplus_{i=0}^{n-1}Ku^i\) 上，按 \(k\) 双线性延拓乘法
\[
(bu^i)(cu^j)=b\sigma^i(c)a^{\lfloor(i+j)/n\rfloor}
u^{\,i+j-n\lfloor(i+j)/n\rfloor}\qquad(b,c\in K,\ 0\leq i,j<n).
\]
则 \(C\) 是维数为 \(n^2\) 的结合含幺 \(k\) 代数。它由子域 \(K\) 和一个可逆元 \(u\) 生成，满足
\[
ub=\sigma(b)u\quad(b\in K),\qquad u^n=a.
\]
更确切地，给定含幺 \(k\) 代数 \(B\)、指定的含幺 \(k\) 代数嵌入 \(\iota:K\hookrightarrow B\)，以及 \(v\in B\) 满足
\[
v\iota(b)=\iota\bigl(\sigma(b)\bigr)v\quad(b\in K),\qquad v^n=\iota(a),
\]
则存在唯一含幺 \(k\) 代数同态 \(\Phi:C\to B\)，使 \(\Phi|_K=\iota\) 且 \(\Phi(u)=v\)。
\end{proposition}
\begin{proof}
比较三个基型元素 \(bu^i,cu^j,du^\ell\) 的乘积，令 \(q_{i,j}=\lfloor(i+j)/n\rfloor\)、\(r_{i,j}=i+j-nq_{i,j}\)。两种结合次序先分别将 \(i+j\) 或 \(j+\ell\) 化到 \(0,\ldots,n-1\)，总进位数满足
\[
q_{i,j}+q_{r_{i,j},\ell}
=\left\lfloor\frac{i+j+\ell}{n}\right\rfloor
=q_{j,\ell}+q_{i,r_{j,\ell}}.
\]
这是将 \(i+j=nq_{i,j}+r_{i,j}\) 或 \(j+\ell=nq_{j,\ell}+r_{j,\ell}\) 代入中间式所得的恒等式。两种乘积因而含有相同的 \(a\) 次幂，剩下的 \(u\) 指数也都是 \(i+j+\ell-n\lfloor(i+j+\ell)/n\rfloor\)。其余系数都为 \(b\sigma^i(c)\sigma^{i+j}(d)\)，因为 \(\sigma^n=1\)，且 \(a\in k\) 被 \(\sigma\) 固定。因此两种乘积相同，双线性再给出全部元素的结合律。\(1\in K\) 为单位，\(k\) 标量位于中心。底向量空间事先已指定为 \(n\) 个 \(K\) 的直和，所以维数确为 \(n^2\)，无需再假设这些形式单项式无关。

关系由定义直接得到，\(u^{-1}=a^{-1}u^{n-1}\)。给定 \(\iota\) 及 \(v\)，唯一可能的映射把 \(\sum_i b_i u^i\) 送到 \(\sum_i\iota(b_i)v^i\)。由 \(v^i\iota(c)=\iota(\sigma^i(c))v^i\) 与 \(v^n=\iota(a)\)，它保持\eqref{b2:eq:cyclic-algebra-product}，所以确为代数同态。
\end{proof}

该代数称为由这些数据给出的\term[xunhuan-daishu]{循环代数}[cyclic algebra]，记为 \((K/k,\sigma,a)\)。上述泛性质是\cref{b2:prop:algebra-relations-universal}中代数展示原则的具体应用：\(K\) 中的加法与乘法已由 \(\iota\) 保持，新生成元的像只需满足所列的两条扭曲关系，就能唯一确定整个代数同态。
\symidx{\((K/k,\sigma,a)\)：循环扩张及参数确定的循环代数}
当 \(n=1\) 时，另约定 \((k/k,\operatorname{id},a)=k\)；这个退化情形不另引入形式生成元。

\subsection{循环代数的中心单性}
\label{b2:subsec:1802:02}

\begin{theorem}\label{b2:thm:cyclic-algebra-central-simple}
设 \(K/k\) 为次数 \(n\geq1\) 的循环 Galois 扩张，\(\operatorname{Gal}(K/k)=\langle\sigma\rangle\)，\(a\in k^\times\)，则循环代数 \(C=(K/k,\sigma,a)\) 是次数 \(n\) 的中心单 \(k\) 代数，且由 \(K\) 分裂。
\end{theorem}
\begin{proof}
当 \(n=1\) 时 \(C=k\)，结论直接成立。以下设 \(n\ge2\)。

要确定中心，先找出哪些元素与子域 \(K\) 交换，再检查它们是否与生成元 \(u\) 交换。设 \(z=\sum_i b_i u^i\) 与 \(K\) 中所有元素交换。比较 \(zc\) 与 \(cz\)，得 \(b_i\bigl(\sigma^i(c)-c\bigr)=0\) 对所有 \(c\in K\) 成立。对 \(i\ne0\)，\(\sigma^i\ne1\)，可以选 \(c\) 使差非零，故 \(b_i=0\)。所以 \(C\) 中 \(K\) 的中心化子就是 \(K\)。若 \(z=b_0\) 再与 \(u\) 交换，则 \(\sigma(b_0)=b_0\)，故 \(z\in k\)。于是中心恰为 \(k\)。

为证单性，需要证明每个非零双边理想都含有 \(1\)。取这样的理想 \(I\)，并在其中选取非零系数个数最少的元素。右乘 \(u^r\)（\(0\le r<n\)）使指数 \(i\) 变为 \((i+r)\bmod n\)，只是循环搬移各分量，并将进位项的系数乘以非零标量 \(a\)。因此可以将任意一个非零系数搬到常数位置，且非零系数个数不变。再左乘这个常数系数在 \(K\) 中的逆，只将所有系数同时乘以非零数，便得到 \(x=1+\sum_{i>0}b_i u^i\in I\)，仍有最少的非零系数。对 \(c\in K\)，交换子 \(xc-cx\in I\) 的常数项为零，其余非零项只能出现在 \(x\) 原有的非零位置，故非零项数严格小于 \(x\)。若它非零，与极小性矛盾；故它对所有 \(c\) 都为零。由前一段的中心化子计算，\(x\in K\)，进而 \(x=1\)。因此 \(I=C\)，代数单。

构造给出 \(\dim_kC=n^2\)，所以 \([K:k]=n=\deg C\)，恰满足\cref{b2:prop:maximal-field-splits-csa}中子域次数等于代数次数的条件，故 \(K\) 分裂 \(C\)。该命题的实际映射为 \(C\otimes_kK\to\operatorname{End}_K(C)\)，\(c\otimes b\) 作用为 \(x\mapsto cxb\)；目标中的空间是原代数 \(C\)，并未换成 \(C_K\)。把它看成右 \(K\) 空间，其维数为 \(n\)，所以目标是 \(M_n(K)\)。
\end{proof}

在次数二、特征不为二时，若 \(K=k(\sqrt d)\)，则 \(u\sqrt d=-\sqrt d\,u\)，\(u^2=a\)。因而循环代数就是四元数代数 \((d,a)_k\)。这里指定自同构是重要数据；较高次数时，换一个生成元需同时跟踪相应参数，不能只把记号中的 \(\sigma\) 删去。

\subsection{范数与分裂条件}
\label{b2:subsec:1802:03}

循环扩张的范数为
\[
\operatorname{N}_{K/k}(b)=\prod_{j=0}^{n-1}\sigma^j(b).
\]
这是第一卷定理18.2.3的共轭元公式在循环 Galois 扩张上的形式。它给出判断 \(C\) 是否为矩阵代数的条件。

\begin{theorem}[循环代数的分裂判据]\label{b2:thm:cyclic-algebra-norm-splitting}
设 \(K/k\) 为有限循环 Galois 扩张，\(\operatorname{Gal}(K/k)=\langle\sigma\rangle\)，\(a\in k^\times\)，则循环代数 \((K/k,\sigma,a)\) 分裂，当且仅当 \(a\in\operatorname{N}_{K/k}(K^\times)\)。
\end{theorem}
\begin{proof}
当 \([K:k]=1\) 时循环代数为 \(k\)，而域范数为恒等映射，结论成立。以下设 \(n=[K:k]\ge2\)。

若 \(a=\operatorname{N}_{K/k}(b)\)，其中 \(b\in K^\times\)，要构造分裂表示，可以让 \(K\) 在自身上左乘，再寻找与扭曲关系相容的算子 \(T\) 来表示 \(u\)。由 \(uc=\sigma(c)u\) 所要求的相容条件是 \(T(cz)=\sigma(c)T(z)\)。\(K\) 是一维 \(K\) 空间，因此任何满足此条件的算子都由 \(T(1)\) 决定，具体为 \(T(z)=\sigma(z)T(1)\)。取 \(T(1)=b\)，就得到 \(k\) 线性算子 \(T(z)=b\sigma(z)\)，并且
\[
T^n(z)=b\sigma(b)\cdots\sigma^{n-1}(b)\sigma^n(z)
=\operatorname{N}_{K/k}(b)z=az.
\]
因此\cref{b2:prop:cyclic-algebra-construction}的泛性质给出含幺 \(k\) 代数同态 \(C\to\operatorname{End}_k(K)\)，也就是\cref{b2:prop:algebra-modules-representations}所说的 \(C\) 在这个空间上的作用。来源单而映射非零，所以单射；两端维数都是 \(n^2\)，故为同构。

反过来，设 \(C\cong\operatorname{End}_k(V)\)，其中 \(\dim_kV=n\)。这个分裂模型给出 \(C\) 的作用，限制到子域 \(K\) 使 \(V\) 成为左 \(K\) 向量空间。比较 \(\dim_kV=[K:k]\dim_KV\)，得到 \(\dim_KV=1\)。选一个 \(K\) 基向量把 \(V\) 识别为 \(K\)。此时 \(u\) 的作用 \(T\) 必须满足 \(T(cz)=\sigma(c)T(z)\)，所以仍由 \(b=T(1)\) 唯一决定，\(T(z)=b\sigma(z)\)。\(u\) 可逆使 \(b\ne0\)，而 \(u^n=a\) 给出 \(\operatorname{N}_{K/k}(b)=a\)：分裂模型中的一维 \(K\) 空间强制产生了所需的范数解。
\end{proof}

当 \(n\ge2\)、\(b\in K^\times\) 时，把生成元改为 \(v=bu\)。每次将 \(u\) 移过下一个 \(b\) 都施加一次 \(\sigma\)，所以
\[
v^n=b\sigma(b)\cdots\sigma^{n-1}(b)u^n
=\operatorname{N}_{K/k}(b)a.
\]
并且 \(vc=\sigma(c)v\) 对每个 \(c\in K\) 成立。新生成元还能恢复 \(u=b^{-1}v\)，故由构造的泛性质得到从参数为 \(\operatorname{N}_{K/k}(b)a\) 的循环代数到 \(C\) 的满同态；两侧维数都为 \(n^2\)，它是同构。因此参数乘上一个范数不改变代数的同构类型。

\ProseParagraphBreak
例如，取 \(\sigma\) 为复共轭，\(\mathbb C/\mathbb R\) 的范数为 \(z\overline z\)，非零范数恰为正实数。故 \((\mathbb C/\mathbb R,\sigma,a)\) 在 \(a>0\) 时分裂，在 \(a<0\) 时同构于 \(\mathbb H\)。前一种情况可以缩放 \(u\) 使平方为一，后一种可以缩放使平方为负一；范数判据同时解释为什么两种情况不能同构。
